% =========================================================
% CHAPTER 5: IDEATION
% =========================================================
\chapter{Ideation}
\label{chap:ideation}

\section{The Ideation Process}
During the ideation phase, multiple architectural paradigms and security models were evaluated to solve the payment fraud challenge. The objective was to discover a solution that maximizes fraud detection accuracy while maintaining low operational latency and a seamless user experience.

\section{Potential Architectural Approaches}
Three distinct solution architectures were evaluated:

\subsection{Approach 1: Static Heuristic Rule Engine}
\begin{itemize}
    \item \textbf{Description}: Hardcoded SQL query filters and threshold checks executed inside database triggers or API controllers (e.g., amount $> 50,000$, transactions $> 5/\text{min}$).
    \item \textbf{Advantages}: Extremely low computation latency ($\sim 1\text{ ms}$), zero training requirements, deterministic execution.
    \item \textbf{Disadvantages}: High false positive rate, inability to detect novel fraud schemes, easily circumvented by splitting transactions.
\end{itemize}

\subsection{Approach 2: Pure Standalone Machine Learning Classifier (Binary Cutoff)}
\begin{itemize}
    \item \textbf{Description}: A standalone machine learning model (e.g., Deep Neural Network or Gradient Boosted Trees) trained to output a binary decision ($\hat{y} \in \{0, 1\}$) at threshold $P(\text{fraud}) \ge 0.50$.
    \item \textbf{Advantages}: Learns complex non-linear feature interactions; adapts to historical data patterns.
    \item \textbf{Disadvantages}: Binary cutoff creates severe user friction for borderline cases; lacks explainability; does not provide intermediate step-up verification.
\end{itemize}

\subsection{Approach 3: Hybrid Adaptive Risk Engine with Explainable AI \& Email OTP (Selected Solution)}
\begin{itemize}
    \item \textbf{Description}: Combines balanced machine learning classification with behavioral heuristic signals into a continuous $0-100$ risk score. Integrates SHAP TreeExplainer for local feature attributions and routes transactions through a 4-tier adaptive policy enforcing pre-settlement Email OTP verification.
    \item \textbf{Advantages}: Maximizes detection accuracy; eliminates binary friction through step-up verification; provides transparent customer and analyst explanations; protects account balances from premature debits.
    \item \textbf{Disadvantages}: Requires integration across ML pipelines, SMTP providers, and transactional state machines.
\end{itemize}

\section{Evaluation and Comparison of Solutions}
To select the optimal architecture, a multi-criteria decision matrix was constructed evaluating the three candidate approaches across critical technical and operational parameters:

\begin{table}[h]
\centering
\caption{Comprehensive Architectural Decision Matrix}
\label{tab:decision_matrix}
\begin{tabularx}{\textwidth}{lcccc}
\toprule
\textbf{Evaluation Criteria} & \textbf{Weight} & \textbf{Approach 1 (Rules)} & \textbf{Approach 2 (Pure ML)} & \textbf{Approach 3 (Hybrid + XAI)} \\
\midrule
\textbf{Detection Accuracy} & 25\% & Low (2/5) & High (4/5) & \textbf{Very High (5/5)} \\
\textbf{False Positive Mitigation} & 20\% & Poor (1/5) & Moderate (3/5) & \textbf{Excellent (5/5)} \\
\textbf{Explainability (XAI)} & 15\% & Moderate (3/5) & Poor (1/5) & \textbf{Outstanding (5/5)} \\
\textbf{User Experience (Friction)} & 15\% & Poor (2/5) & Moderate (3/5) & \textbf{Excellent (5/5)} \\
\textbf{Settlement Protection} & 15\% & Poor (2/5) & Moderate (3/5) & \textbf{Very High (5/5)} \\
\textbf{Implementation Feasibility}& 10\% & High (5/5) & High (4/5) & \textbf{High (4/5)} \\
\midrule
\textbf{Weighted Score} & 100\% & \textbf{2.15 / 5.0} & \textbf{3.25 / 5.0} & \textbf{4.85 / 5.0} \\
\bottomrule
\end{tabularx}
\end{table}

\section{Selected Architecture Justification}
Approach 3 was unanimously selected as the foundational architecture for \textbf{FraudShield AI}. The hybrid paradigm provides an optimal synergy:
\begin{enumerate}[label=\alph*.]
    \item \textbf{High-Fidelity ML Inference}: Scikit-Learn tuned Random Forest classification delivers near-perfect precision ($1.0000$) and high recall ($0.9970$) on holdout evaluations.
    \item \textbf{Contextual Behavioral Signals}: Account-draining ratios, transaction velocity, and diurnal anomalies adjust the baseline ML probability, preventing edge-case bypasses.
    \item \textbf{Dual Explainability via SHAP}: Solves the black-box dilemma by decomposing predictions into additive feature attributions for both customers and SOC analysts.
    \item \textbf{Adaptive Step-Up OTP Settlement Security}: By placing borderline transactions in an \texttt{OTP\_REQUIRED} state and locking settlement until verification, customer funds are safeguarded without causing hard rejections.
\end{enumerate}
